\begin{lemma}[Intrinsic minimum-ancestor description of the joint owner cut]
\label{lem:acf-joint-owner-minimum-ancestors}
Let $\mathcal L$ be a FINITE
distributive lattice, presented as
$\operatorname{Down}(P)$ by
Birkhoff's theorem. Suppose
$P$ has $r\ge2$ distinct
minimal points $p_1,\dots,p_r$.
For $q\in P$ write
\[
    \operatorname{Anc}(q)=
         \{\,i\in[r]:p_i\le q\,\}.
\]
This set is NONEMPTY for every
$q\in P$. Let the maximal
pseudocomplements of its exact
atoms be
\[
 C_i=P\setminus\uparrow p_i
       =\{q:i\notin\operatorname{Anc}(q)\}.
\]
Put $J_{ij}=C_i\vee C_j$,
$H_i=\bigwedge_{j\ne i}J_{ij}$
and $H_*=\bigwedge_{i<j}J_{ij}$.
Under the literal downset
representation these have
the EXACT descriptions
\[
\begin{aligned}
 J_{ij}
    &=\{q:\{i,j\}\not\subseteq
                     \operatorname{Anc}(q)\},\\
 H_i
    &=\{q:i\notin\operatorname{Anc}(q)
                 \text{ OR }
                 \operatorname{Anc}(q)=\{i\}\},\\
 H_*
    &=\{q:|\operatorname{Anc}(q)|=1\}.
\end{aligned}
\tag{MO}\label{eq:acf-joint-owner-minimum-ancestors}
\]
Consequently
\[
 H_*>0,\qquad
 \bigwedge_iH_i=H_*,
 \qquad H_i\vee H_j=J_{ij}
                     \quad(i\ne j).
\tag{MOR}\label{eq:acf-joint-owner-relations}
\]
The first inequality is STRICT
and uses no original
amalgamation: every atomic
minimal point $p_i$ belongs
to $H_*$. In the standing
finite-exact-normalised AP+CCO
source, take $C_i=\widehat C_i$,
so the joint owner quotient
$M/H_*$ is therefore an ACTUAL
STRICT exact quotient (lower
diversity rank) and the
pairwise joins are precisely
the common owner quotients
used by finite exact CRT.
This proof is pure finite
distributive-lattice algebra;
it does not assume independent
coordinate factors or Boolean
original orbit structure.
\end{lemma}
\begin{proof}
An element $q$ belongs to
the downset $C_i$ iff
$q\not\ge p_i$, equivalently
$i\notin\operatorname{Anc}(q)$.
Joins and meets of downsets
are LITERAL unions and
intersections. Therefore
$q\in C_i\vee C_j$ iff
$i\notin\operatorname{Anc}(q)$
or $j\notin\operatorname{Anc}(q)$,
giving the first line.

Likewise
$q\in H_i$ iff for EVERY
$j\ne i$, either
$i\notin\operatorname{Anc}(q)$
or $j\notin\operatorname{Anc}(q)$.
If $i$ is absent this is
automatic; if $i$ is
present, it demands that
NO other minimal ancestor
be present. This is the
second line.

Finally $q\in H_*$ iff
NO two distinct minimal
ancestors belong to
$\operatorname{Anc}(q)$.
Every point of a finite
poset lies above some
minimal element, so
$|\operatorname{Anc}(q)|\ge1$.
This proves the third line.

The equality
$\bigwedge_iH_i=H_*$
is immediate from the
definitions or the ancestor
criterion. For $i\ne j$,
$q\notin H_i\vee H_j$
iff BOTH i and j belong
to $\operatorname{Anc}(q)$,
which is exactly the
condition $q\notin J_{ij}$.
Thus $H_i\vee H_j=J_{ij}$.
Each minimal point $p_i$
has $\operatorname{Anc}(p_i)=\{i\}$,
hence belongs to $H_*$.
So $H_*>0$.

In the original source,
\cref{lem:acf-hat-pseudocomplement}
identifies the actual exact
$\widehat C_i$ with
$P\setminus\uparrow p_i$,
and the exact lattice is
distributive. A nonidentity
exact quotient collapses at
least one original diversity
atomic orbital to its
quotient diagonal and hence
strictly decreases original
diversity rank. The owner
cut $H_*$ contains every
atom $E_i$, so is
nonidentity, establishing
the stated strict quotient.
\end{proof}

\begin{corollary}[Maximal cores give admissible simultaneous factors]
\label{cor:acf-hat-factors}
Assume $r\ge2$ and use $\widehat C_i$ as above.  Put
\[
 J_{ij}=\widehat C_i\vee\widehat C_j,\qquad
 H_i=\bigcap_{j\ne i}J_{ij},\qquad
 H_*=\bigcap_{i<j}J_{ij},\qquad
 \Omega=M/H_*.
\]
Then $H_*>\Id$, $H_i\vee H_j=J_{ij}$ and
$\bigcap_iH_i=H_*$, with their exact
minimum-ancestor normal form supplied
by \cref{lem:acf-joint-owner-minimum-ancestors}.
The actual owner source $\Omega$
is therefore a strict exact quotient.
Every $Q_i=M/\widehat C_i$ has a global full-support
strong closed safe palette as in
Lemma~\ref{lem:acf-quotient-safety}; and every original
unsafe clique has at most one projection not contained
in a single finite $F_i$ packet.  Also
\[
 \operatorname{acl}_M(A)=
 \bigcap_i\pi_{\widehat C_i}^{-1}
          \bigl(F_i[\pi_{\widehat C_i}(A)]\bigr)
 =\bigcap_iR_i[A].
\]
\end{corollary}
\begin{proof}
For each $k$, every $J_{ij}$ contains $E_k$:
if $k=i$ use the other core, and if $k=j$ use the first;
otherwise either suffices.  Thus
$\Id<\bigvee_kE_k\subseteq H_*$.  Since
$\widehat C_i\subseteq H_i$, the inclusion
$J_{ij}\subseteq H_i\vee H_j$ follows, while the
reverse inclusion follows from $H_i,H_j\subseteq J_{ij}$.
The meet equality follows by distributing the finite
intersection of the same pairwise joins.

For fixed $i$, put $K=\bigcap_{j\ne i}\widehat C_j$,
$J=\widehat C_i\vee K$.  The two complementary
exact quotients have identity meet.  $J$ contains all
minimal original $E_j$, so it is harmless in $M$ by
Lemma~\ref{lem:acf-harmless-owners}.  In $Q_i$
its kernel contains the sole minimal infinite direction;
in the other factor $M/K$, Lemma~\ref{lem:acf-hat-normalised}
identifies the remaining minimal directions, all of
which are contained in this kernel.  Therefore the
owner parameter is harmless in BOTH quotient factors.
Apply Lemma~\ref{lem:acf-quotient-safety}.

The unsafe-clique concentration is the previous
pairwise $R_i,R_j$ theorem and is independent of the
choice of exact core: for every $j$ other than the
possible exceptional index, the whole clique lies
in a single $R_j$-class, which projects to a finite
$R_j/\widehat C_j$ packet.  Lastly, the $\widehat C_i$
meet in $\Id$, and joint exact quotient homogeneity
gives finite-meet acl reconstruction.  Use the unary
formula in each $Q_i$, then take full inverse images;
these are exactly the displayed $R_i[A]$.
\end{proof}
